\documentclass[11pt,a4paper]{article}
\usepackage[top=1.7cm,bottom=1.5cm,left=1.5cm,right=1.5cm,headsep=8pt,footskip=16pt]{geometry}
\usepackage{enumitem}
\usepackage{amssymb}
\usepackage{array}
\usepackage[table]{xcolor}
\usepackage{fancyhdr}
\usepackage{lastpage}
\usepackage{multicol}
\usepackage{polyglossia}
\setdefaultlanguage[calendar=gregorian, numerals=maghrib]{arabic}
\setotherlanguage{english}
\IfFontExistsTF{Amiri}{%
  \newfontfamily\arabicfont[Script=Arabic, Scale=1.08]{Amiri}
}{%
  \newfontfamily\arabicfont[Script=Arabic, Scale=1.08, Path=../fonts/, Extension=.ttf,
    UprightFont=*-Regular, BoldFont=*-Bold, ItalicFont=*-Italic,
    BoldItalicFont=*-BoldItalic]{Amiri}
}
\IfFontExistsTF{AlBattar}{%
  \newfontfamily\khatt[Script=Arabic, Scale=1.15]{AlBattar}
}{%
  \newfontfamily\khatt[Script=Arabic, Scale=1.15, Path=../fonts/,
    BoldFont=ae_AlBattar.ttf]{ae_AlBattar.ttf}
}
\renewcommand{\headrulewidth}{0.8pt}
\renewcommand{\footrulewidth}{0.8pt}
\fancyhf{}
\rhead{\khatt المدرسة العليا للأساتذة ببوزريعة}
\lhead{\khatt ملحقة بومرداس}
\lfoot{\khatt الأستاذ: د. بن عبد الرزاق يوسف}
\cfoot{\thepage/\pageref{LastPage}}
\pagestyle{fancy}
\parindent 0pt
\setlength{\columnsep}{0.6cm}
\setlength{\columnseprule}{0.4pt}
\setlist[enumerate,1]{label=\textbf{\arabic*.}, leftmargin=*, itemsep=5pt, topsep=2pt}
\begin{document}

\linespread{1.08}\selectfont
\begin{center}
{\LARGE\khatt\textbf{الاختبار الرابع في الذكاء الاصطناعي}}\\[4pt]
{\normalsize وضعيات إشكالية ــ أسئلة متعددة الاختيارات ــ الفصول من 1 إلى 6}
\end{center}
\vspace{-2pt}\rule{\textwidth}{0.6pt}\\[2pt]
\textbf{التعليمة:} يتضمن الاختبار خمس وضعيات إشكالية. اقرأ كل وضعية بعناية، ثم اختر التحليل أو التصرف الأنسب
من بين الاختيارات الخمسة؛ لكل وضعية إجابة صحيحة واحدة فقط.
اكتب \textbf{حرف} الإجابة الصحيحة في جدول الإجابات في آخر الاختبار.
لكل وضعية 4 نقاط (المجموع 20 نقطة).
\vspace{4pt}

\noindent
{\large\khatt\textbf{الوضعية 1}}\hfill{\small(الفصل 1)}\par\nopagebreak\vspace{3pt}
\setlength{\fboxsep}{7pt}\colorbox{gray!10}{\parbox{\dimexpr\linewidth-14pt}{يعدّ أستاذ لتلاميذ السنة الثالثة ابتدائي قصة مصورة عن الأمانة من ستة مشاهد، بنص مضبوط بالشكل وقراءة مسجلة. لاحظ في تجربة أولى أن ملامح الشخصية الرئيسية تتغير من صورة إلى أخرى، وأن النص يتضمن كلمات فوق مستوى التلاميذ. ويريد أن ينشر القصة في مكتبة المدرسة. ما الخطة الأنسب لإنتاج القصة؟}}\par\nopagebreak
\begin{itemize}[leftmargin=1.0cm, labelsep=6pt, itemsep=3pt, topsep=4pt, parsep=0pt]
  \item[\textbf{أ)}] يطلب من الأداة النصية كتابة القصة ووصف الصور وتحديد طريقة القراءة في طلب واحد طويل، لأن جمع المهام يضمن الاتساق بين النص والصور والصوت، ثم يعتمد المخرجات دون مراجعة لأنها صادرة عن طلب واحد.
  \item[\textbf{ب)}] يولّد الصور أولًا بوصف مختصر لكل مشهد، ثم يطلب من أداة النص كتابة قصة تناسب الصور، لأن البدء بالصور يحل مشكلة تغيّر الملامح تلقائيًا، ويجعل النص في مستوى التلاميذ دون تحديد ذلك في الطلب.
  \item[\textbf{ج)}] يحدد مستوى التلاميذ في طلب النص، ويقبل تغيّر ملامح الشخصية لأنه لا يمكن التحكم فيه أبدًا في أدوات توليد الصور، ثم يولّد القراءة من النص دون ضبطه لأن الأداة الصوتية تضبط النص بنفسها.
  \item[\textbf{د)}] يصف الشخصية وصفًا ثابتًا في كل طلب ويراجع اتساق الصور، لكنه يُبقي النص كما هو، لأن الكلمات التي تفوق مستوى التلاميذ تثري رصيدهم، والقصص الموجهة للأطفال ينبغي أن تتحدى مستواهم دائمًا.
  \item[\textbf{هـ)}] يحدد في طلب النص مستوى التلاميذ ومفردات مناسبة ويراجع الضبط، ثم يصف الشخصية وصفًا ثابتًا مفصلًا يكرره في طلب كل مشهد ويراجع اتساق الصور، ثم يولّد القراءة من النص المضبوط ويستمع إليها كاملة.
\end{itemize}
\par\vspace{6pt}
\noindent
{\large\khatt\textbf{الوضعية 2}}\hfill{\small(الفصل 2)}\par\nopagebreak\vspace{3pt}
\setlength{\fboxsep}{7pt}\colorbox{gray!10}{\parbox{\dimexpr\linewidth-14pt}{طرحت أستاذتان في يوم واحد السؤال نفسه حرفيًا على الأداة نفسها: «اقترح عنوانًا لدرس عن الوصف»، فحصلت كل منهما على عنوان مختلف. ثم سألت إحداهما عن سنة ميلاد شاعر مغمور، فحصلت في محاولتين متتاليتين على سنتين مختلفتين. واستغربتا ذلك، فعرضتاه على زميل لهما في الإعلام الآلي. أي تفسير يجمع الحالتين بدقة؟}}\par\nopagebreak
\begin{itemize}[leftmargin=1.0cm, labelsep=6pt, itemsep=3pt, topsep=4pt, parsep=0pt]
  \item[\textbf{أ)}] النموذج يحفظ إجابة ثابتة لكل سؤال، لكنه يكيّف الإجابة حسب سجل كل مستخدم، فاختلاف الإجابات في الحالتين دليل على التخصيص، وتُعتمد أي من السنتين لأن كلتيهما صحيحة بالنسبة إلى من سألها.
  \item[\textbf{ب)}] الاختلاف في الحالتين عطل تقني عابر، لأن النماذج مصممة لتعطي الإجابة نفسها لكل سؤال متطابق؛ والحل إعادة السؤال حتى تتكرر إجابة واحدة مرتين، فتكون هي الصحيحة في العناوين والتواريخ معًا.
  \item[\textbf{ج)}] النموذج يولّد الكلمة تلو الكلمة باختيار احتمالي، وهذا يفسّر اختلاف العناوين فقط؛ أما اختلاف سنة الميلاد فسببه أن النموذج حدّث معلوماته من الإنترنت بين المحاولتين، فتُعتمد السنة الثانية لأنها الأحدث.
  \item[\textbf{د)}] النموذج يولّد الكلمة تلو الكلمة باختيار احتمالي، فتتغير الإجابة بين المحاولات؛ وهذا التنوع مفيد في اقتراح العناوين، لكنه في المعلومات الدقيقة علامة على عدم ثبات المعرفة، فيلزم التحقق من السنة في مصدر موثوق.
  \item[\textbf{هـ)}] النموذج يولّد الكلمة تلو الكلمة باختيار احتمالي، وهذا التنوع غير مرغوب في الحالتين، والحل أن يُضاف إلى كل طلب «أجب بالإجابة الصحيحة الوحيدة»، لأن هذه العبارة تلغي الاختيار الاحتمالي وتضمن صحة المعلومات، وهي طريقة مجرَّبة.
\end{itemize}
\par\vspace{6pt}
\noindent
{\large\khatt\textbf{الوضعية 3}}\hfill{\small(الفصل 4)}\par\nopagebreak\vspace{3pt}
\setlength{\fboxsep}{7pt}\colorbox{gray!10}{\parbox{\dimexpr\linewidth-14pt}{ولّد طالب قصيدة بأداة ذكاء اصطناعي، ثم عدّل نحو ثلث أبياتها وأصلح وزنها، وأراد المشاركة بها في مسابقة جامعية يشترط نظامها «نصوصًا من إبداع المتسابق»، دون أن يذكر شيئًا عن الذكاء الاصطناعي. وقد طلب رأي أستاذه قبل إرسال المشاركة. وهو يعتقد أن تعديلاته كافية. ما الموقف الأسلم أخلاقيًا؟}}\par\nopagebreak
\begin{itemize}[leftmargin=1.0cm, labelsep=6pt, itemsep=3pt, topsep=4pt, parsep=0pt]
  \item[\textbf{أ)}] يشارك بالقصيدة دون إفصاح، لأن تعديل ثلث الأبيات وإصلاح الوزن يجعل النص من إبداعه، والنظام لم يمنع الذكاء الاصطناعي صراحة، وما لم يُمنع صراحة مباح في المسابقات الجامعية والأدبية.
  \item[\textbf{ب)}] يشارك بالقصيدة ويفصح في آخرها عن استعمال الأداة دون سؤال اللجنة، لأن الإفصاح وحده يجعل المشاركة مطابقة لشرط الإبداع، ولو كان النص في أصله مولَّدًا وكانت التعديلات محدودة.
  \item[\textbf{ج)}] سكوت النظام لا يعني الإباحة، والنص في أصله مولَّد، فالأسلم أن يسأل اللجنة ويفصح لها عن طريقة الإنتاج قبل المشاركة، فإن لم تقبل شارك بنص من إبداعه، مستفيدًا من الأداة في التدرّب لا في الإنتاج.
  \item[\textbf{د)}] يعدّل نصف الأبيات بدل الثلث ثم يشارك دون إفصاح، لأن تجاوز نسبة النصف في التعديل هو الحد المعتمد عالميًا لاعتبار النص المولَّد من إبداع صاحبه، فلا تحتاج اللجنة إلى معرفة طريقة الإنتاج، كما تفعل بعض المسابقات الدولية.
  \item[\textbf{هـ)}] يمتنع عن أي مشاركة أدبية في المستقبل، لأن مجرد تجربته للأداة في كتابة الشعر يُفقده صفة المبدع نهائيًا، حتى لو كتب لاحقًا نصوصًا من تأليفه الخالص دون الاستعانة بأي أداة، احترامًا لنزاهة المسابقات.
\end{itemize}
\par\vspace{6pt}
\noindent
{\large\khatt\textbf{الوضعية 4}}\hfill{\small(الفصل 5)}\par\nopagebreak\vspace{3pt}
\setlength{\fboxsep}{7pt}\colorbox{gray!10}{\parbox{\dimexpr\linewidth-14pt}{يعدّ أستاذ تقريرًا لمجلس الأساتذة عن نتائج شهادة التعليم المتوسط لهذه السنة، فسأل أداة ذكاء اصطناعي غير متصلة بالبحث في الإنترنت، فقدّمت نسبًا دقيقة بفواصل عشرية، وأكدت عند سؤالها أنها «من المصادر الرسمية». والتقرير سيُعرض غدًا. وزميله يقترح الاكتفاء بالأرقام لضيق الوقت، ما دامت تبدو دقيقة ومفصلة. ما الموقف الصحيح؟}}\par\nopagebreak
\begin{itemize}[leftmargin=1.0cm, labelsep=6pt, itemsep=3pt, topsep=4pt, parsep=0pt]
  \item[\textbf{أ)}] لا يعتمد النسب، لأن أداة غير متصلة بالإنترنت لا تعرف ما بعد تاريخ انتهاء بيانات تدريبها، ودقة الأرقام وتأكيد الأداة لا يثبتان شيئًا، فيرجع إلى الموقع الرسمي لوزارة التربية أو إلى معطيات مديرية التربية.
  \item[\textbf{ب)}] يعتمد النسب بعد أن يطلب من الأداة رابط المصدر الرسمي، فإن قدّمت رابطًا يبدو رسميًا أدرجه في التقرير دون فتحه، لأن وجود رابط بصيغة رسمية دليل كافٍ على أن الأرقام مأخوذة من الوزارة، توفيرًا للوقت قبل الاجتماع.
  \item[\textbf{ج)}] يعتمد النسب بعد تقريبها إلى أقرب عدد صحيح، لأن الخطأ في الأرقام المولَّدة يكون عادة في الفواصل العشرية فقط، أما الأعداد الصحيحة فتعكس الاتجاه العام الذي تعرفه الأداة من نتائج السنوات الماضية.
  \item[\textbf{د)}] لا يعتمد النسب، ويطلب من الأداة تقديرًا جديدًا مبنيًا على نتائج السنوات الماضية، ويقدّمه في التقرير على أنه النتيجة الرسمية لهذه السنة، لأن التقدير المبني على اتجاه السنوات السابقة قريب عادة من الواقع.
  \item[\textbf{هـ)}] يعتمد النسب لأن الأداة أكدت أنها من المصادر الرسمية، والأدوات لا تؤكد مصدرًا إلا إذا اطلعت عليه، ويذكر في التقرير أن الأرقام مأخوذة من الأداة حتى يتحمل مطوّروها المسؤولية عند الخطأ.
\end{itemize}
\par\vspace{6pt}
\noindent
{\large\khatt\textbf{الوضعية 5}}\hfill{\small(الفصل 6)}\par\nopagebreak\vspace{3pt}
\setlength{\fboxsep}{7pt}\colorbox{gray!10}{\parbox{\dimexpr\linewidth-14pt}{رفعت أستاذة إلى أداة متعددة الأنماط صورة مخطط بياني من الكتاب المدرسي عن تطور عدد المطالعين خلال عشر سنوات، وطلبت شرحه لتلاميذها. جاء الشرح سليم اللغة حسن التنظيم، لكن الأداة قرأت قيمتين خطأً، فاستنتجت تراجعًا غير موجود في المخطط. وتريد بناء نشاط تقويمي على هذا الشرح. ما الدرس المنهجي الأدق؟}}\par\nopagebreak
\begin{itemize}[leftmargin=1.0cm, labelsep=6pt, itemsep=3pt, topsep=4pt, parsep=0pt]
  \item[\textbf{أ)}] قراءة الأداة للأرقام من الصور قد تخطئ، لكن الاستنتاج العام عن التراجع يبقى صحيحًا، لأن الأداة تفهم اتجاه المنحنى حتى لو أخطأت في القيم الجزئية، فتعتمد الاستنتاج وتصحّح القيمتين فقط، وهذا يكفي في الشرح للتلاميذ.
  \item[\textbf{ب)}] قراءة الأداة للأرقام من الصور قد تخطئ، وسلامة اللغة وحسن التنظيم لا يدلّان على صحة القراءة؛ فتقابل كل قيمة بالمخطط قبل اعتماد الاستنتاج، ويمكنها تزويد الأداة بالقيم مكتوبة لتعتمد عليها بدل قراءتها من الصورة.
  \item[\textbf{ج)}] الخطأ سببه أن الصورة ملونة، والأدوات متعددة الأنماط تقرأ المخططات بالأبيض والأسود بدقة تامة، فتحوّل الصورة إلى رمادية وتعيد الطلب، وتعتمد الشرح الجديد دون مقابلته بالمخطط الأصلي.
  \item[\textbf{د)}] الخطأ دليل على أن مخطط الكتاب المدرسي غير واضح، لأن الأدوات متعددة الأنماط تقرأ المخططات الواضحة دون أخطاء، فتبلّغ الجهة المعدّة للكتاب، وتعتمد مؤقتًا شرح الأداة مع تنبيه التلاميذ إلى ذلك.
  \item[\textbf{هـ)}] الأدوات متعددة الأنماط لا تفهم المخططات إطلاقًا، وكل ما قالته مختلق من وحي العنوان، فتستبعد استعمالها في أي مهمة فيها صور، وتكتفي بالأدوات النصية التي لا تخطئ في الأرقام عندما تُكتب لها، حتى تتجنب الأخطاء في المستقبل.
\end{itemize}
\par\vspace{6pt}
\begin{center}
{\khatt\textbf{جدول الإجابات}}\\[3pt]
\renewcommand{\arraystretch}{1.6}
\begin{tabular}{|>{\centering\arraybackslash}m{1.9cm}|*{5}{>{\centering\arraybackslash}m{1.2cm}|}}
\hline
\textbf{رقم السؤال} & \textbf{1} & \textbf{2} & \textbf{3} & \textbf{4} & \textbf{5} \\ \hline
\textbf{الإجابة} &  &  &  &  &  \\[8pt] \hline
\end{tabular}

\vspace{8pt}
\renewcommand{\arraystretch}{2.2}
\begin{tabular}{|>{\centering\arraybackslash}m{7.5cm}|>{\centering\arraybackslash}m{3.2cm}|>{\centering\arraybackslash}m{3.2cm}|}
\hline
\textbf{الاسم واللقب} & \textbf{الفوج} & \textbf{العلامة} \\ \hline
 & & \hfill /20 \\ \hline
\end{tabular}
\end{center}
\end{document}
